\documentclass[11pt]{article}
\usepackage[a4paper,margin=26mm]{geometry}
\usepackage{amsmath,amssymb,amsthm,booktabs,array}
\usepackage[T1]{fontenc}
\usepackage{lmodern}
\usepackage[colorlinks=true,linkcolor=blue,citecolor=blue,urlcolor=blue]{hyperref}
\hypersetup{pdftitle={Graffeo--Gubbiotti, Remark 7.7 (arXiv:2301.09129v2): Nambu conjugacy and the twelve-plane configuration},pdfauthor={Haobo Ma; Wenlin Zhang}}
\newtheorem{theorem}{Theorem}
\newtheorem{corollary}[theorem]{Corollary}
\newcommand{\C}{\mathbb C}
\newcommand{\R}{\mathbb R}
\newcommand{\adj}{\operatorname{adj}}
\newcommand{\diag}{\operatorname{diag}}
\newcommand{\Cr}{\mathcal C}
\title{Graffeo--Gubbiotti, Remark 7.7\newline
\large arXiv:2301.09129v2, pp.\ 35--37\newline
\Large Nambu conjugacy and the twelve-plane configuration}
\author{Haobo Ma\thanks{ChronoAI Pte Ltd, Singapore.}\qquad Wenlin Zhang\thanks{National University of Singapore; ChronoAI Pte Ltd, Singapore.}}
\date{November 3, 2026}
\begin{document}
\maketitle
\begin{abstract}
\setlength{\emergencystretch}{2em}
We resolve Graffeo--Gubbiotti's Remark~7.7
(arXiv:2301.09129v2, pp.\ 35--37) for quadratic Nambu maps with a
positive-definite first form and simple generalized spectrum.
Simultaneous congruence reduces the vector field to an Euler top, and an explicit
four-dimensional linear change conjugates its Kahan--Hirota--Kimura map to
$(2I_4-\mathbf1\mathbf1^T)\circ\Cr$, where $\Cr$ is the standard cubic Cremona map.
The argument also allows a symmetric second Hamiltonian and every nonzero step.
It gives the square factors in (7.15) and an explicit answer to the
twelve-plane incidence question following (7.23): each of the eight base
points lies on six planes, each plane contains four points, and the conjugator
sends the points to the two normal-form tetrahedra.
We also treat regular complex pencils with three distinct projective roots
and determine the linear degree growth when two real generalized eigenvalues coincide.
\end{abstract}
\noindent\small Licensed under Creative Commons Attribution 4.0 International
(CC BY 4.0).\normalsize

\section{Statement and normalization}
Graffeo--Gubbiotti, Remark~7.7, arXiv:2301.09129v2, pp.\ 35--37
\cite{GG}, asks for projective conjugacy of the quadratic Nambu
Kahan--Hirota--Kimura map to type (A)-(i), and for the arrangement of eight
points on twelve planes. Both follow from the explicit conjugator below.
The authors' earlier version~1, Remark~7.7, pp.\ 34--35, already describes the
Nambu--Euler class identification through singularity patterns; version~2
isolates the conjugacy conjecture and the incidence question.
Their Lemma~2.6 and Corollary~2.9 prove the Euler-top factorization and conjugacy.
Celledoni et al.\ \cite[Proposition~5]{Celledoni} establish the Nambu
integrals underlying the net (7.16). Related constructions of three-dimensional
QRT maps use pairs of pencils of quadrics \cite[Sections~3--5]{AS}; here a
constant simultaneous congruence and a cubic identity give the conjugator.

Let $A,B$ be real symmetric $3\times3$ matrices, set
$H_1(x)=x^TAx$, $H_2(x)=x^TBx$, and put
\[
 f(x)=\nabla H_1(x)\times\nabla H_2(x)=4(Ax\times Bx),\qquad
 L(x)v=Ax\times Bv+Av\times Bx.
\]
Thus $Df(x)=4L(x)$. For $X=(x,w)\in\C^4$ and real $h\ne0$, define the cubic lift
\begin{equation}\label{eq:lift}
 M(X)=wI_3-2hL(x),\qquad
 F_h(X)=\bigl(\adj(M(X))x,\det M(X)\bigr).
\end{equation}
On $w=1$, this represents $(I_3-\frac h2 Df(x))^{-1}x$, the Kahan map:
its polarized equation is $x'-x=\frac h2 Df(x)x'$ because $f$ is homogeneous
quadratic. This also agrees with \cite[(2.2)--(2.3)]{GG}.
Write
\[
 \Cr(Y)=(Y_2Y_3Y_4,Y_1Y_3Y_4,Y_1Y_2Y_4,Y_1Y_2Y_3),\quad
 G=2I_4-\mathbf1\mathbf1^T,
\]
where $\mathbf1=(1,1,1,1)^T$; then $G^2=4I_4$ and $\det G=-16$.

\begin{theorem}\label{thm:main}
Suppose $A$ is positive definite, $B$ is symmetric, and
\begin{equation}\label{eq:rank}
 \operatorname{rank}(B-tA)\ge2\quad\text{for every }t\in\R.
\end{equation}
For every real $h\ne0$, there is $U\in\operatorname{GL}_4(\C)$ such that
\begin{equation}\label{eq:normal}
 U F_h(X)=-\tfrac18 G\Cr(UX)\qquad(X\in\C^4).
\end{equation}
Consequently $[F_h]$ is projectively conjugate to $[G\Cr]$.
In particular the positive-definite case of Remark 7.7 with simple generalized
spectrum is of type (A)-(i).
\end{theorem}

\section{An explicit conjugator}
Choose an orthogonal eigenbasis $Q$ of the real symmetric matrix
$A^{-1/2}BA^{-1/2}$ and set $P=A^{-1/2}Q$. Then
\[
 P^TAP=I_3,\qquad P^TBP=\diag(\mu_1,\mu_2,\mu_3),\qquad d=\det P\ne0.
\]
Congruence preserves rank, so (\ref{eq:rank}) is equivalent to pairwise
distinct $\mu_i$: the rank of $\diag(\mu_i-t)$ is the number of nonzero
diagonal entries. Put
\begin{equation}\label{eq:a}
 a_1=\frac4d(\mu_3-\mu_2),\qquad
 a_2=\frac4d(\mu_1-\mu_3),\qquad
 a_3=\frac4d(\mu_2-\mu_1).
\end{equation}
All three coefficients are nonzero. The identity
\begin{equation}\label{eq:cross}
 (P^{-T}u)\times(P^{-T}v)=d^{-1}P(u\times v)
\end{equation}
follows by pairing both sides with $P^{-T}z$ and using the determinant formula
for the scalar triple product. Since $AP=P^{-T}$ and
$BP=P^{-T}\diag(\mu_i)$, substitution of $x=Py$ gives
\[
 P^{-1}f(Py)=(a_1y_2y_3,a_2y_1y_3,a_3y_1y_2)=:e(y).
\]
Differentiation gives $Df(Py)P=PDe(y)$. With $R=\diag(P,1)$ and
\[
 E_h(y,w)=\bigl(\adj(wI_3-\tfrac h2 De(y))y,
                    \det(wI_3-\tfrac h2 De(y))\bigr),
\]
conjugation of the matrix in (\ref{eq:lift}) therefore gives the polynomial
identity
\begin{equation}\label{eq:cov}
 F_h(RZ)=R E_h(Z).
\end{equation}
Here determinant and adjugate commute with similarity, including at singular
matrices; the latter identity follows first for invertible matrices and then
as a polynomial identity.

Choose complex square roots $\alpha^2=a_1$, $\beta^2=a_2$, $\gamma^2=a_3$, and set
\begin{equation}\label{eq:T}
 T=\begin{pmatrix}
 -\beta\gamma h&-\alpha\gamma h&\alpha\beta h&-2\\
 \beta\gamma h&-\alpha\gamma h&-\alpha\beta h&-2\\
 -\beta\gamma h&\alpha\gamma h&-\alpha\beta h&-2\\
 \beta\gamma h&\alpha\gamma h&\alpha\beta h&-2
 \end{pmatrix}.
\end{equation}
Its determinant is $32\alpha^2\beta^2\gamma^2h^3\ne0$. The Euler identity is
\begin{equation}\label{eq:euler}
 T E_h(Z)=-\tfrac18 G\Cr(TZ).
\end{equation}
For completeness, the following quartic calculation proves this identity
with the normalization used here. Set
\[
 q=(r,s,t,v)=(h\beta\gamma y_1,h\alpha\gamma y_2,h\alpha\beta y_3,-2w),
 \quad
 J=\begin{pmatrix}-1&-1&1&1\\1&-1&-1&1\\-1&1&-1&1\\1&1&1&1\end{pmatrix}.
\]
Then $TZ=Jq$, $J^TJ=4I_4$, and
$GJ=2J D_0$ for $D_0=\diag(1,1,1,-1)$.
Expansion of the $3\times3$ adjugate shows that the corresponding scaled
output of $E_h$ is $q'=\frac14(p_r,p_s,p_t,p_v)$, where
\begin{align*}
 p_r&=r(v^2+s^2+t^2-r^2)-2stv,\\
 p_s&=s(v^2+r^2+t^2-s^2)-2rtv,\\
 p_t&=t(v^2+r^2+s^2-t^2)-2rsv,\\
 p_v&=v^3-v(r^2+s^2+t^2)+2rst.
\end{align*}
To check the grouping of the cubic terms, write $\pi(Y)=Y_1Y_2Y_3Y_4$.
Multiplying the four linear forms in $Jq$ gives
\[
 \pi(Jq)=r^4+s^4+t^4+v^4
 -2(r^2s^2+r^2t^2+s^2t^2+r^2v^2+s^2v^2+t^2v^2)+8rstv.
\]
Taking its gradient yields
$J^T\Cr(Jq)=-16D_0q'$. Hence
$J^TG\Cr(Jq)=-32q'$, and multiplication by $J/4$ gives
$Jq'=-G\Cr(Jq)/8$, proving (\ref{eq:euler}).
Finally, take $U=TR^{-1}$ and apply (\ref{eq:euler}) to
$Z=R^{-1}X$ in (\ref{eq:cov}). This proves Theorem~\ref{thm:main}, including
the homogeneous identity at every $X$, not just on the affine domain.

\section{Consequences at the posing source}
\begin{corollary}\label{cor:planes}
Under the hypotheses of Theorem~\ref{thm:main}, the degree-eight factors in
\cite[(7.15)]{GG} are squares of products of four distinct linear forms.
The indeterminacy loci of the map and its inverse are each six lines forming
the edges of a tetrahedron. The degree of its $n$th iterate is $2n^2+1$.
\end{corollary}
\begin{proof}
The matrices above obey $T_{-h}=-GT_h/2$, and therefore $U_{-h}=-GU_h/2$.
Writing $U=U_h$, (\ref{eq:normal}) at both signs of $h$ gives
\[
 F_h=-\tfrac18U^{-1}G\Cr U,\qquad
 F_{-h}=-\tfrac1{32}U^{-1}\Cr G U.
\]
Since $\Cr(\Cr(Y))=\pi(Y)^2Y$ and $\Cr(cY)=c^3\Cr(Y)$, composition yields
\begin{equation}\label{eq:compositions}
 F_{-h}(F_h(X))=\frac{\pi(UX)^2}{256}X,\qquad
 F_h(F_{-h}(X))=\frac{\pi(GUX)^2}{65536}X.
\end{equation}
Invertibility of $U$ and $GU$ makes the four factors in each product distinct.
The base locus of $\Cr$ is the union of the six lines $Y_i=Y_j=0$
($i<j$). Thus the two base loci are the inverse images of these lines under
$U$ and $GU$, respectively; this transports \cite[Remark 2.8]{GG}.
Finally, degree is preserved by projective conjugacy, so the type (A)
degree formula in \cite[Proposition 5.3]{GG} applies.
\end{proof}

\begin{corollary}\label{cor:eight}
The net in \cite[(7.16)]{GG}, with
$H_3(x)=x^TA\adj(B)Ax$, has eight distinct base points.
\end{corollary}
\begin{proof}
The base equations are $H_1=H_2=w^2+4h^2H_3=0$.
In $x=Py$ coordinates, adjugate congruence gives
\[
 H_1=\sum_i y_i^2,\quad H_2=\sum_i\mu_i y_i^2,\quad
 H_3=d^{-2}\sum_i\mu_j\mu_k y_i^2\quad(\{i,j,k\}=\{1,2,3\}).
\]
Put $b=(\mu_2-\mu_3,\mu_3-\mu_1,\mu_1-\mu_2)$ and
$\Delta=(\mu_1-\mu_2)(\mu_1-\mu_3)(\mu_2-\mu_3)$.
The first two equations force $(y_1^2,y_2^2,y_3^2)=\lambda b$; the identity
$\sum_i\mu_j\mu_k b_i=\Delta$ then gives
$w^2=-4h^2d^{-2}\lambda\Delta$.
Here $\lambda\ne0$, since $\lambda=0$ would force all four homogeneous
coordinates to vanish. All $b_i$ and $\Delta$ are nonzero. After projective
rescaling to $\lambda=1$, there are sixteen choices of square roots, identified
in pairs by simultaneous sign reversal. These are eight distinct points.
They are reduced: on the locus of nonzero coordinates, coordinatewise squaring
has invertible differential, and the three equations in the squares are
independent linear equations.
\end{proof}

\begin{corollary}\label{cor:incidence}
Under the hypotheses of Theorem~\ref{thm:main}, the net (7.16) has exactly
six reducible members, each a pair of distinct planes. In suitable coordinates
$v=(v_1,v_2,v_3,v_4)$, they are
\[
 (v_i-v_j)(v_i+v_j)=0\qquad(1\le i<j\le4),
\]
and the base points are
$p_{\epsilon}=[\epsilon_1:\epsilon_2:\epsilon_3:1]$,
$\epsilon_i\in\{1,-1\}$.
Each point lies on six of the twelve distinct planes, and each plane contains
four of the eight points. More precisely, with $\epsilon_4=1$,
\[
 p_{\epsilon}\in\{v_i-\sigma v_j=0\}
 \quad\Longleftrightarrow\quad \sigma=\epsilon_i\epsilon_j.
\]
The conjugator $U$ sends the base-point set to
$\{[e_i],[Ge_i]:1\le i\le4\}$, where $e_i$ are the coordinate vectors.
\end{corollary}
\begin{proof}
Let $c=4h^2/d^2$ and use $b,\Delta$ from Corollary~\ref{cor:eight}.
Choose $s_i^2=b_i$ and $s_4^2=-c\Delta$, and put
$v_i=y_i/s_i$ and $v_4=w/s_4$. The net becomes the full
three-dimensional vector space of diagonal quadrics
$\sum_{i=1}^4 k_i v_i^2$ with $\sum_i k_i=0$:
the three generators are independent and vanish at all sign vectors.
A nonzero diagonal quadric is reducible over $\C$ precisely when it has
rank at most two. Rank one is excluded by $\sum k_i=0$; rank two forces
exactly two nonzero coefficients $k_i=-k_j$. This proves the six-member
description and distinctness of all twelve planes. The displayed incidence
rule gives four points per plane and one plane per coordinate pair at each point.

For explicit members in the $y,w$ coordinates, for each
$\{i,j,k\}=\{1,2,3\}$ take
\[
 (\mu_j-\mu_i)y_j^2+(\mu_k-\mu_i)y_k^2,
 \qquad w^2+c(\mu_i-\mu_j)(\mu_i-\mu_k)y_i^2.
\]
Their net parameters are respectively
$[-\mu_i:1:0]$ and $[-c\mu_i(\mu_j+\mu_k):c\mu_i:1]$.
These provide the six members in \cite[(7.20)--(7.23)]{GG}
without choosing a parameter chart that can vanish at diagonal forms.

Finally $a_i=-4b_i/d$ and $b_1b_2b_3=-\Delta$ show that the
scaled coordinates $q$ in Section~2 satisfy
$q_1^2=q_2^2=q_3^2=q_4^2\ne0$ at each base point.
The eight projective sign vectors split into four even-parity and four
odd-parity classes. Multiplication by $J$ sends the former to $[e_i]$
and the latter to $[Ge_i]$, as follows by multiplying its four rows
by each sign vector. Since $UX=Jq$, this proves the last assertion and
answers the incidence question on p.\ 37 of \cite{GG}.
\end{proof}

\section{Extensions and collisions}
\begin{corollary}\label{cor:complex}
For complex symmetric $A,B$, suppose the binary cubic
$\det(sA+tB)$ has three distinct roots in $\mathbb P^1(\C)$.
Then, for every $h\in\C^*$, the Nambu Kahan map is projectively conjugate
to $[G\Cr]$.
\end{corollary}
\begin{proof}
First suppose $A$ is invertible and $A^{-1}B$ has three distinct eigenvalues.
The eigenvectors form a basis and are orthogonal for the nondegenerate
bilinear form $u^TAv$: if their eigenvalues are $\mu_i\ne\mu_j$, symmetry
gives $(\mu_i-\mu_j)u_i^T A u_j=0$.
None has zero self-pairing, since it would then be orthogonal to the whole
basis and lie in the radical of $A$. Scaling them gives
$P^TAP=I$, $P^TBP=\diag(\mu_i)$ over $\C$.
All identities in Section~2 apply unchanged, proving the assertion.
In general choose an invertible member $A'=pA+qB$ and an independent
$B'=rA+sB$, with $\delta=ps-qr\ne0$.
Such an $A'$ exists because the binary determinant is nonzero.
The three projective roots remain distinct after this parameter change, so
$A'^{-1}B'$ has simple spectrum. The Nambu vector field of $(A',B')$ is
$\delta$ times that of $(A,B)$. Its Kahan map at step $h/\delta$ is
therefore exactly the original map at step $h$; the first case applies.
\end{proof}

\begin{corollary}\label{cor:collision}
Let $A>0$ and $B=B^T$ be real, and let $h\ne0$. For $n\ge0$, the degrees
of the iterates are $2n^2+1$ for three distinct generalized eigenvalues,
$2n+1$ for exactly two distinct generalized eigenvalues, and $1$ when all
three coincide. All these maps have dynamical degree one and algebraic
entropy zero.
\end{corollary}
\begin{proof}
The first case is Corollary~\ref{cor:planes}. In the second, order the
eigenvalues so that $\mu_2=\mu_3\ne\mu_1$. Simultaneous congruence still
gives the Euler field $(0,a y_1y_3,-a y_1y_2)$ with
$a=4(\mu_1-\mu_2)/d\ne0$.
Put $z_+=y_2+\mathrm i y_3$, $z_-=y_2-\mathrm i y_3$ and
$L_\pm=w\pm \mathrm i h a y_1/2$.
The polarized equations give $y_1'=y_1$ on $w=1$, and
$z_+'=(L_-/L_+)z_+$, $z_-'=(L_+/L_-)z_-$.
Thus, in coordinates $[y_1:z_+:z_-:w]$, the $n$th iterate has lift
\[
 \bigl(y_1(L_+L_-)^n,\ z_+L_-^{2n},\ z_-L_+^{2n},\ w(L_+L_-)^n\bigr).
\]
For $n\ge1$ these degree-$2n+1$ polynomials have no common factor:
$L_+,L_-$ are distinct, and the second and third entries omit opposite
factors; their other possible common factors are excluded by the first
and fourth entries. For $n=0$ the map is the identity.
If all eigenvalues coincide, $B=\mu A$, so the vector field vanishes.
Finally all three degree laws have $\lim_{n\to\infty}(\deg F^n)^{1/n}=1$.
\end{proof}

This covers the semisimple collisions in the positive-definite family
mentioned in \cite[Section~8]{GG}. Singular pencils, nonsemisimple
collisions, and the separate map $\theta$ in (8.5)--(8.7) remain different
questions.

\begin{center}\small
\begin{tabular}{@{}p{0.25\linewidth}p{0.67\linewidth}@{}}
\toprule
Source item in \cite{GG} & Coverage\\\midrule
Remark 7.7 & Proved: Theorem~\ref{thm:main}; extended to regular simple complex pencils in Corollary~\ref{cor:complex}.\\
(7.15) & Proved: exact square factors in (\ref{eq:compositions}).\\
(7.16) & Proved: eight reduced base points (Corollary~\ref{cor:eight}); the eight-point assertion is already in \cite{GG}, p.\ 36.\\
(7.20)--(7.23); p.\ 37 & Proved: six reducible quadrics and twelve-plane incidence, Corollary~\ref{cor:incidence}; the six-member assertion is already in \cite{GG}.\\
Lemma 2.6; Corollary 2.9; Remark 2.8; Proposition 5.3 & Prior proof cited: Euler factorization and conjugacy, tetrahedra, quadratic degree law; transported here.\\
(7.14); \cite[Prop.\ 5]{Celledoni} & Prior proof cited: modified Nambu integrals.\\
Section 8: collisions & Proved for $A>0$, $B=B^T$: Corollary~\ref{cor:collision}.\\
Section 8: other degenerations; (8.5)--(8.7) & Not covered: singular pencils, nonsemisimple collisions, $\theta$ and its singularity pattern.\\
\bottomrule
\end{tabular}
\end{center}

\paragraph{Statement and computation coverage.}
The pinned Lean source
\href{https://github.com/the-omega-institute/trureturing/blob/86172888e706177b6c1d135cb12aec3b6d2cdce8/D5/S3/FluidDynamics/Solitons/NambuKahanCremonaDecomposition.lean}{\texttt{.lean} file}
contains the declaration
\begin{quote}\small\ttfamily
D5.S3.FluidDynamics.Solitons.\\
NambuKahanCremonaDecomposition.result
\end{quote}
Its proposition states that for real positive-definite $A,B$ satisfying
(\ref{eq:rank}) and $h>0$, the lift (\ref{eq:lift}) equals
$L_2\Cr(L_1X)$ with $L_1,L_2\in\operatorname{GL}_4(\C)$.
The auxiliary declaration \texttt{euler\_top\_identity}
proves (\ref{eq:euler}) for complex Euler coefficients with specified
square roots; \texttt{covariance} proves (\ref{eq:cov}) for a given
simultaneous congruence. Its \texttt{assembly} constructs the conjugator
internally; the public \texttt{result} asserts the two-projectivity factorization.
The explicit conjugacy for symmetric $B$ and $h\ne0$, and all
Corollaries~\ref{cor:planes}--\ref{cor:collision}, have the written proofs above;
the complex-pencil and collision extensions are statements of those written proofs.
The accompanying \texttt{verify.py} checks the universal Euler cubic,
determinant, cross-product, quartic, reversal and composition identities by
exact polynomial arithmetic, as well as an exact nonorthogonal example.
It also enumerates all eight sign points and twelve planes, verifies their
48 incidences and their images under $J$, and checks the six net members
and the collision formula through the first four iterates.

\begin{thebibliography}{9}
\bibitem{GG} M. Graffeo and G. Gubbiotti,
\emph{Growth and integrability of some birational maps in dimension three},
Ann. Henri Poincar\'e \textbf{25} (2024), 1733--1793.
\href{https://doi.org/10.1007/s00023-023-01339-5}{doi:10.1007/s00023-023-01339-5}.
Posing version: \href{https://arxiv.org/abs/2301.09129v2}{arXiv:2301.09129v2}
(June 3, 2023), Remark 7.7, pp. 35--37.
Earlier class-identification discussion:
\href{https://arxiv.org/abs/2301.09129v1}{version 1} (January 22, 2023),
Remark 7.7, pp. 34--35.
\bibitem{Celledoni} E. Celledoni, R. I. McLachlan, D. I. McLaren,
B. Owren and G. R. W. Quispel,
\emph{Integrability properties of Kahan's method},
J. Phys. A: Math. Theor. \textbf{47} (2014), 365202.
\href{https://doi.org/10.1088/1751-8113/47/36/365202}{doi:10.1088/1751-8113/47/36/365202};
\href{https://arxiv.org/abs/1405.3740}{arXiv:1405.3740}, Proposition 5.
\bibitem{AS} J. Alonso and Y. B. Suris,
\emph{Invariant volume form for 3D QRT maps},
\href{https://arxiv.org/abs/2510.11468v1}{arXiv:2510.11468v1}
(October 13, 2025), Sections 3--5.
\end{thebibliography}
\end{document}
